\chapter{Discussion}\label{chap:6}

The study links the geometry available to partners with the actions through which they develop a construction. An arrangement can look plausible, become the basis for further work and be accepted even when its proportions differ from the target. Parts of that arrangement also have histories: one person builds them, another continues them, and either may later revise them. Following those histories connects the geometric value of an overhead view to the organisation of collaborative work.

\section{How the overhead view's geometric advantage enters collaboration}\label{sec:6-1}

Tangram makes local connections and the overall outline relevant to the same task. A piece must fit its neighbours, and the complete arrangement must match the target. Perspective projection from an oblique viewpoint changes the apparent lengths and angles of a planar construction. A perpendicular orthographic view preserves its relative shape and scale. This geometric property matters during collaboration because the developing shape is what partners use when they choose the next piece, assess a proposal and decide whether to continue changing the result.

G06 connects this property to an observed problem. Both partners helped rebuild the candle to include the unused medium triangle, and the seven-piece result was accepted by the pair and the researcher. The researcher later noticed that the accepted outline had a wider lower body than the target, and reconstruction confirmed that mismatch. The task account establishes that it was missed at the time. Their combination gives a concrete example of coordinated construction reaching an accepted result whose geometry still needed correction. [\hyperlink{src:D14}{D14}, \hyperlink{src:D16}{D16}]

The three first-person cases also show different responses to a plausible shape. G01 rejected a replacement when its edges failed to align. G05 continued after a partner questioned the chair's proportions. G06 accepted its final candle. The researcher's records identify this category in three of six first-person pairs and none of the six shared-overhead pairs. The comparison locates the problem within the implemented viewing conditions; the episodes explain how it entered the work. Agreement about an ending and the geometric validity of that ending are separate achievements.

Chen et al. compared a world-in-miniature interface and an interactive 2D map in collaborative searching and object manipulation \cite{chen2024}. Their discussion identifies a need to examine finer placement accuracy beyond placing objects within a designated work area. Tangram brings that issue into a planar construction where relative spans and edges determine whether a candidate matches the target. The present study also follows how partners form and revise the candidate, connecting spatial presentation with the sequence that produces the judged result.

The geometric advantage must be considered alongside control. In this prototype, the shared overhead image replaced the first-person scene while tracked controllers remained in the shared 3D workspace. G05 and G06 requested or provided help with difficult manipulations. Questionnaire responses likewise separated explaining a placement from moving or rotating a piece. A clear overall view can therefore coexist with difficulty carrying out the next local action. This combination makes the viewing arrangement and its manipulation mapping joint design concerns. [\hyperlink{src:D11}{D11}, \hyperlink{src:D15}{D15}]

\section{Expressing ideas through manipulation and taking turns to advance the task}\label{sec:6-2}

The partners' placements gave each other something concrete to inspect and continue. In G01's chicken, a demonstrated alternative for the back was followed by additions to the head. In G06's physical-overhead goose, P12 added a small triangle to an existing neck arrangement and P11 continued that addition with the medium triangle. Each action used what was already present and changed the possibilities for the next action. The sequence shows how a local construction can function as a proposal within the joint task. [\hyperlink{src:D01}{D01}, \hyperlink{src:D12}{D12}] These sequences develop a shared problem state through action, consistent with Roschelle and Teasley's account of proposals being presented and taken up through talk and manipulation \cite{roschelle1995}.

Dillenbourg describes how the distinction between cooperation and collaboration depends partly on the grain at which a division of labour is examined \cite{dillenbourg1999}. The Tangram cases make that grain visible at the level of parts and successive placements. A participant could develop a part independently for a short period, while the part remained dependent on the partner's preceding and subsequent work. The interpretation is therefore a changing organisation of connected actions: temporary division of work, continuation and joint revision can all occur within one attempt.

The G06 action trace provides a detailed example of that change. Initial work on the body and top developed into repeated manipulation of the same body pieces by both participants. Its paired pose sequence shows what those operations produced. G03 adds a different path: P05 brought two small triangles together, separated them and was followed by P06 operating on those pieces. Tracking operators together with configurations distinguishes a handover that continues an existing arrangement from one that changes the direction of construction. [\hyperlink{src:D14}{D14}, \hyperlink{src:D06}{D06}]

Workspace awareness includes following who is active and what is happening in the shared workspace \cite{gutwin2002}. The Tangram analysis follows the consequence of that activity for a developing arrangement. Yang et al. separated object interaction from meaningful changes of grouping in collaborative image organisation \cite{yang2022}. The same distinction matters here: a grab can adjust alignment, test a fit, expose a reference or change a piece's role. Reading the operation within its sequence explains its contribution to the task more clearly than the number of operations alone.

Irlitti et al. distinguished closely engaged collaboration, where one person manipulated objects while others discussed the task, from episodes where one person took over and others contributed little \cite{irlitti2024}. That distinction helps interpret the Tangram cases. A period of single-person manipulation can belong to joint work when it responds to the partner's preceding construction or provides a next step that the partner takes up. The evidence lies in those connections, so the analysis follows responses and changing arrangements alongside the distribution of control.

Manipulation can also help a participant develop an idea while acting. Kirsh and Maglio distinguish actions that advance a physical task from actions that make information easier to discover or evaluate \cite{kirsh1994}. Their account offers a way to understand why testing and rearranging pieces belongs within problem solving, even when a move is later reversed. In these cases, the observable evidence is the attempted relation and what follows it; the records support describing exploration without assigning a hidden cognitive purpose to every movement.

Rasch et al. studied undo in deliberately divided and collaborative VR tasks, then identified constantly changing collaboration arrangements as an issue for further work \cite{rasch2024}. The freely organised Tangram episodes show such changes within a task. The contribution lies in tracing how a local arrangement connects successive initiatives and revisions. This gives a concrete account of how work becomes joint at a fine scale, including moments when one participant is the only person manipulating a piece.

\section{Retaining a partner's partial arrangement and revising one's own}\label{sec:6-3}

A partial construction can remain available to the partner and later be reopened by its original builder. G01's correct goose head and neck survived the partner's tidying before its builder changed the relation into an incorrect one. G03's self-revision had a different consequence: dismantling a small-triangle arrangement made those pieces available for the developing solution. These cases show why revising a part and improving the whole are distinct. A change can resolve a conflict between parts or remove a useful local result. [\hyperlink{src:D03}{D03}, \hyperlink{src:D06}{D06}]

Retaining a construction also lets earlier work enter a later step. In G06's physical-overhead goose, successive additions developed from the existing neck and led to a correct outline with a different internal partition from the reference. The square and medium triangle occupied different regions, each paired with a small triangle. The outcome matters because the task allowed alternative internal solutions. Preservation supported a viable route through the problem even though that route differed from the reference construction. [\hyperlink{src:D12}{D12}]

The G06 interview gives one participant's interpretation of this behaviour: the partner's leaving placed pieces in position was taken as a sign of understanding. Yang et al. used a more specific criterion for respect in their image-organisation study: asking for permission before modifying another person's grouping \cite{yang2022}. These accounts concern different evidence. Here, retained geometry establishes that a construction remained available, while the interview explains what that action meant to a participant. Agreement or respect should be attributed through such contextual evidence. [\hyperlink{src:D13}{D13}]

An enacted alternative could prompt direct partner intervention. G01's chicken-back demonstration and G03's horse-head replacement occurred before the complete solution was established. In G05, the partner questioned the changed role of the parallelogram and later restored it to the backrest. These cases accompany the preservation and self-revision cases. The relevant comparison is whether an alternative was proposed or put into action in the recorded sequence, and what the partner did next. An absence of visible intervention provides no account of alternatives a participant might have considered privately. [\hyperlink{src:D01}{D01}, \hyperlink{src:D05}{D05}, \hyperlink{src:D10}{D10}]

G06's candle joins preservation to joint rebuilding. The square and the later red large-triangle placement persisted in the pose record while both participants continued to adjust other pieces. This establishes a stable part of the evolving construction around which further action occurred. The eventual result still had the measured proportion mismatch. Examining both histories explains how coordinated continuation and an incorrect outcome can coexist: the pair advanced a shared arrangement, while some of its geometric relations survived the ongoing changes. [\hyperlink{src:D14}{D14}, \hyperlink{src:D16}{D16}]

\section{Design implications for shared-overhead VR}\label{sec:6-4}

Shared-overhead VR should make overall geometric relations easy to inspect during construction. G06's request for another viewpoint shows a participant treating viewing as a resource for resolving the task. A design with accessible switching could let partners move between local manipulation and inspection of the whole arrangement. Piece identities, the target and orientation cues should remain recognisable across that transition so that the construction can be continued. Its value should be evaluated through the checking and subsequent actions that the view supports. [\hyperlink{src:D14}{D14}]

Manipulation needs comparable attention. Assistance in G05 and G06 shows a practical route through a blocked action: another participant carries out an adjustment, then work continues. Interfaces should make a piece's current control state clear and allow brief transfers without forcing a permanent director--operator arrangement. The relevant outcome is whether partners can complete the requested adjustment and resume their work. The prototype's direction mapping and reach difficulties should be addressed alongside visual clarity, because both affect who can carry out an available next step. [\hyperlink{src:D11}{D11}, \hyperlink{src:D15}{D15}]

Construction history offers a further way to support revision. The cases suggest retaining access to earlier local arrangements so that partners can compare or recover a useful relation after changing it. Rasch et al. found that the scope of undo in multi-user VR affects both recovery and disturbance to a partner's work \cite{rasch2024}. Their results motivate careful choice of what a restoration changes. For Tangram, an inspectable earlier state could support discussion before a shared arrangement is restored; this proposal remains to be evaluated in use.

Evaluation should connect these features to episodes of work. For a view change, the question is which relation becomes inspectable and what the partners do with it. For assistance, it is whether control passes successfully and the blocked placement is completed. For history, it is which earlier arrangement is recovered and how later additions are affected. These questions preserve the connection between interface behaviour, changing geometry and the partners' next actions that made the present cases informative.

\section{Study scope and further evaluation}\label{sec:6-5}

The findings describe six familiar pairs working through four conditions in a fixed order. Target patterns and familiarisation varied; G01's first recording used an earlier target presentation, while the recorded G02--G06 sessions used the revised version. The three-of-six versus zero-of-six comparison concerns the defined near-complete shape problem in the researcher's records. It is separate from a uniform geometric inspection of every ending and from an estimate of viewpoint's independent effect. A subsequent comparison can stabilise the controls and target presentation, then balance patterns and order.

The analysis combines sources with different strengths. Researcher process reports locate proposals, responses and visible actions. Logs identify accepted control and recorded geometry; selected reconstructions check shape problems that the researcher identified during review. Interview paraphrases provide participants' retrospective meanings. Missing audio limits the attribution of some spoken turns, while the retained source references keep those boundaries visible. The G06 reconstruction applies the existing outline requirement, preserving both the original acceptance and the subsequent correction.

The selected revision cases establish sequences, including instances of self-revision and direct partner intervention. They support comparing how a partial solution was treated and what happened to it. Establishing their relative frequency would require consistent episode definitions applied across all recorded processes. The current contribution is the linked account: who first formed an arrangement, how it remained available, who changed it and how its geometry developed. This level of analysis makes useful local work and its possible loss visible within the broader task.


The findings give a more specific answer to the initial expectations. The overview's value is visible in the planar relationships available for checking a construction, while the overhead cases also show the importance of usable control and local help. The construction sequences explain how partners temporarily divide, continue and revise work across conditions. They leave open the earlier question of whether viewpoint changes personal agency or ownership: the present analysis traces how partial arrangements are treated, while that broader experience was not measured as a separate construct. [\hyperlink{src:D32}{D32}, \hyperlink{src:D31}{D31}]
